%=====================================================================
% Representation-Theoretic Uniqueness of $(2,3)$
% Converted from standalone paper: rep_theoretic_uniqueness.tex
%=====================================================================
\section{Representation-Theoretic Uniqueness of $(2,3)$}
\label{ch:rep_uniqueness}

This chapter strengthens the dimensional argument of Chapters~\ref{ch:whyC}--\ref{ch:whyd5}. We prove that the conclusion does not depend on $d = 5$ at all: for \emph{any} dimension $N \geq 5$, the unique physically non-trivial content is $(n_S, n_T) = (3, 2)$, and its coupling constant is the mathematical constant $\agut = 6/(25\pi^2)$. The Standard Model gauge group $\mathrm{SU}(3)\times\mathrm{SU}(2)\times\mathrm{U}(1)$ is a theorem of mathematics, not a property of a particular universe.

%=====================================================================
\subsection{Background: Atoms, Swap, and $\CC$}
%=====================================================================

\subsubsection{The unique substrate}

The DRLT axiom --- ``$N$ points exist with pairwise relations $G_{ij}$ taking values in an algebra $\mathbb{K}$'' --- requires four conditions for a universe with geometry and forces:

\begin{enumerate}
\item[\textbf{R1.}] \textbf{Magnitude} (geometry): $G_{ij}$ has a well-defined $|G_{ij}|$. Requires $\mathbb{K}$ normed.
\item[\textbf{R2.}] \textbf{Phase} (gauge forces): $\arg(G_{ij})$ is continuous. Requires unit-norm elements to form a connected Lie group.
\item[\textbf{R3.}] \textbf{Determinant} (gravity): $\det(G)$ is a single real number. Requires $\mathbb{K}$ commutative.
\item[\textbf{R4.}] \textbf{Winding} (charge quantization): holonomy admits discrete winding numbers. Requires $\pi_1(\text{phase group}) \cong \mathbb{Z}$.
\end{enumerate}

By the Frobenius classification \cite{Frobenius, Hurwitz}, the
finite-dimensional associative division algebras over $\RR$ are
$\{\RR, \CC, \mathbb{H}\}$ (octonions excluded by
non-associativity; see \cite{Baez} for a modern treatment). $\RR$ fails R2 ($\{+1,-1\}$ is discrete). $\mathbb{H}$ fails R3 (non-commutative, determinant ambiguous). $\CC$ satisfies all four, with $\pi_1(S^1) = \mathbb{Z}$ (R4).

\begin{remark}
$G_{ij} = \langle\psi_i|\psi_j\rangle$ is the unique sesquilinear form on $\CC^d$ satisfying R1--R3. Positive semi-definiteness ($\mathbf{c}^\dagger G\,\mathbf{c} = \|\Psi^\dagger\mathbf{c}\|^2 \geq 0$) is a theorem, not an assumption. The Born rule $W = |G|^2/d$ is the unique minimal-degree polynomial $\CC \to \RR_{\geq 0}$ satisfying real, symmetric, non-negative, and faithful.
\end{remark}

\subsubsection{Atomic dimensions}

\begin{lemma}
The atomic dimensions are exactly $\{2, 3\}$.
\end{lemma}

\begin{proof}
$n = 1$: $\text{SU}(1) = \{1\}$, trivial. $n \geq 4$: $n = 2 + (n-2)$ with $n-2 \geq 2$, decomposable. Only $n = 2$ and $n = 3$ admit no decomposition into parts $\geq 2$.
\end{proof}

\subsubsection{Swap annihilation}

If a decomposition contains two blocks of the same size $n$, the gauge group contains $\text{SU}(n)_1 \times \text{SU}(n)_2$, which admits the outer automorphism $\sigma: (g_1, g_2) \mapsto (g_2, g_1)$. This swap maps every chiral representation to its mirror, forcing all surviving representations to be vector-like. Vector-like representations cannot produce CP violation. Therefore, repeated atomic pairs are physically trivial: present in the eigenvalue spectrum, but incapable of generating distinguishable physics.

%=====================================================================
\subsection{Numerical Investigation: No Spectral Gap at Rank 5}
%=====================================================================

\subsubsection{Rank-5 regions do not spontaneously emerge}

Random normalized vectors in $\CC^d$ ($d = 20, 30, 50$) for $N = 100$--$500$ vertices. The Gram matrix $G_{ij} = \langle\psi_i|\psi_j\rangle$ is computed, and subsets analyzed for effective rank (number of eigenvalues capturing 99\% of the trace). Result: effective rank $\approx d$ everywhere. No rank-5 regions emerge spontaneously.

\subsubsection{SVD truncation produces $(1, 12, 12)$}

The only operation producing rank 5 is SVD truncation: projecting onto the top-5 singular subspace. The Binet--Cauchy decomposition then yields SSS $\approx 1$, SST $\approx 12$, STT $\approx 12$, total $= 25 = d^2$. However, SVD truncation is externally imposed, not a physical mechanism.

\subsubsection{Swap annihilation: gap location}

Example ($d=15$): blocks $= [2, 2, 3, 3, 2, 3]$. After swap-symmetrizing paired blocks: the spectral gap appears at rank $10$ (independent dimensions $= 15 - 2 - 3$), \textbf{not} at rank 5.

Key finding: swap annihilation makes paired blocks \emph{identical}, not zero. Eigenvalues persist; only distinguishability is lost.

%=====================================================================
\subsection{Eigenvalue Structure: Soft Boundary and Hard Wall}
%=====================================================================

\subsubsection{Two qualitatively different boundaries}

The post-swap eigenvalue spectrum has two boundaries:

\begin{center}
\begin{tabular}{ccc}
$\lambda_1 \cdots \lambda_5$ & $\lambda_6 \cdots \lambda_{10}$ & $\lambda_{11} = 0$ \\
chiral $(2+3)$ & trivial (pairs) & annihilated \\
\emph{soft boundary} & & \emph{hard wall} \\
\end{tabular}
\end{center}

\textbf{Soft boundary} ($\lambda_5$ vs $\lambda_6$): statistical. Closes as $1/\sqrt{N}$. The chiral and trivial eigenvalues merge into a continuous spectrum at large $N$.

\textbf{Hard wall} ($\lambda_{d_\text{indep}}$ vs $\lambda_{d_\text{indep}+1} = 0$): algebraic. Exact at all $N$. Paired blocks become identical, reducing rank permanently.

\subsubsection{Convergence analysis}

The gap $1 - \lambda_6/\lambda_5 \approx 0.8/\sqrt{N}$ passes through $\agut \approx 0.0243$ at $N_\text{cross} \approx 1500$. This may correspond to the GUT symmetry-breaking scale: below $N_\text{cross}$, force separation is spectral; above it, only algebraic.

\subsubsection{The Standard Model is invisible to spectroscopy}

No spectral measurement can distinguish chiral from trivial dimensions at large $N$. The Standard Model gauge structure is encoded in the \textbf{representation theory} of the surviving atomic pair, not in eigenvalue magnitudes.

%=====================================================================
\subsection{The Uniqueness Theorem}
%=====================================================================

\begin{theorem}[Uniqueness of $(2,3)$]
For any integer $N \geq 5$, the unique atomic decomposition of $\CC^N$ supporting (a) chiral fermions, (b) CP violation, (c) connected spacetime, and (d) charge quantization, is $(n_S, n_T) = (3, 2)$, yielding $\text{SU}(3)\times\text{SU}(2)\times\text{U}(1)$.
\end{theorem}

\begin{proof}
\textbf{Step 1:} Atoms are $\{2, 3\}$. (The additive atoms lemma, Chapter~\ref{ch:whyd5}.)

\textbf{Step 2:} Any $N$ decomposes into atoms ($N = 2a + 3b$).

\textbf{Step 3:} Repeated atoms are physically trivial (swap $\to$ vector-like $\to$ no CP violation).

\textbf{Step 4:} At most one copy of each atom survives:

\begin{center}
\begin{tabular}{lccccl}
\hline
Surviving & Gauge group & Space? & Time? & Glue? & Physics? \\
\hline
$\{\}$ & trivial & --- & --- & --- & Dead \\
$\{2\}$ & $\text{SU}(2)$ & No & Yes & --- & Dead \\
$\{3\}$ & $\text{SU}(3)$ & Yes & No & --- & Dead \\
$\{2, 3\}$ & $\text{SU}(3)\!\times\!\text{SU}(2)\!\times\!\text{U}(1)$ & Yes & Yes & Yes & Alive \\
\hline
\end{tabular}
\end{center}

\textbf{Step 5: $\{2\}$ and $\{3\}$ alone are dead.}

\emph{Topologically:} $\pi_1(\text{SU}(2)) = \pi_1(\text{SU}(3)) = 0$ --- both simply connected. No winding numbers, no charge quantization, no conserved particle number.

\emph{Physically:} $\{3\}$ alone means $n_T = 0$: no temporal sector, no unitary evolution, no time. $\{2\}$ alone means $n_S = 0$: no spatial sector, no matter, no extent. Neither can support physics.

\textbf{Step 6: U(1) is the unique and necessary glue.}

The relative phase between $\CC^3$ and $\CC^2$ is $\text{U}(1)$, with $\pi_1(\text{U}(1)) = \pi_1(S^1) = \mathbb{Z}$. This provides winding numbers for charge quantization and the topological connection between sectors. Without U(1), the two sectors decouple: $\CC^3$ has no time, $\CC^2$ has no space, and each dies independently. The tracelessness condition $n_S \cdot y_\text{color} + n_T \cdot y_\text{weak} = 0$ enforces $y_\text{color}/y_\text{weak} = -2/3$, ensuring exact charge conservation.

\textbf{Step 7:} Independence from $N$. Steps 1--6 use only the algebra of $\text{SU}(2)$, $\text{SU}(3)$, $\text{U}(1)$. Whether $N = 5$ or $N = 10^{60}$, the physically non-trivial content is $(2, 3)$.
\end{proof}

\begin{corollary}
The question ``why $d = 5$?'' dissolves. The physical content is $(2, 3)$ regardless of $d$. The number $5 = 2 + 3$ is a consequence, not a cause.
\end{corollary}

%=====================================================================
\subsection{Three Independent Paths to $\agut$}
%=====================================================================

\subsubsection{Path 1: Simplex geometry (Binet--Cauchy)}

The exterior algebra $\Lambda^3(\CC^5)$ decomposes under the $(3,2)$ split into SSS ($k=0$), SST ($k=1$), STT ($k=2$) sectors. With the lattice speed of light $c = n_T/d_T \div n_S/d_S = 2$, the $c$-weighted channel count is:
\begin{equation}
\underbrace{1}_{\text{SSS}} + \underbrace{12}_{\text{SST}} + \underbrace{12}_{\text{STT}} = 25 = d^2.
\end{equation}

The gauge multiplicity of each force follows a single principle: \emph{pure hinges get the adjoint representation ($n^2-1$), mixed hinges get the fundamental ($n$)}:
\begin{align}
\text{Strong (SSS, pure } \CC^3\text{):} \quad g_3 &= n_S^2 - 1 = 8, \\
\text{Weak (STT, mixed):} \quad g_2 &= n_T = 2, \\
\text{EM (SST, mixed):} \quad g_1 &= n_S = 3.
\end{align}

Each channel propagates via the solid-angle propagator $D(n) = 1/n^{n_S-1} = 1/n^2$ (not a lattice Green's function, but the geometric solid angle in $n_S = 3$ dimensions). The total interaction capacity:
\begin{equation}
S(\infty) = \sum_{n=1}^{\infty} \frac{1}{n^2} = \zeta(2) = \frac{\pi^2}{6}.
\end{equation}
Therefore: $1/\agut = d^2 \times \zeta(2) = 25\pi^2/6$, giving $\agut = 6/(25\pi^2)$.

\subsubsection{Path 2: Random matrix theory (GUE)}

$G = \Psi\Psi^\dagger$ with $\Psi \in \CC^{N \times d}$ belongs
to the GUE ($\beta = 2$) \cite{Mehta}, forced by
$\mathbb{K} = \CC$. The two-point cluster function $Y_2(r) = (\sin(\pi r)/(\pi r))^2$ gives $R_2(r) \approx 2\zeta(2)r^2$ at short distance:
\begin{equation}
\agut = 1/(d^2 \cdot \zeta(2)) = 6/(25\pi^2).
\end{equation}

\subsubsection{Path 3: Number theory (coprimality)}

$P(\gcd(a,b) = 1) = 1/\zeta(2) = 6/\pi^2$
(Euler--Ces\`aro; see \cite{Hardy-Wright}, \S18.5). Among $d^2 = 25$ channels, the coprime fraction is $6/\pi^2$:
\begin{equation}
\agut = P(\text{coprime})/d^2 = 6/(25\pi^2).
\end{equation}

\begin{theorem}[GUE-Coupling Correspondence]
$\agut = 1/(d^2\zeta(2))$ is the normalized short-distance level repulsion coefficient of the GUE sine kernel, partitioned over $d^2$ independent channels.
\end{theorem}

$\agut = 6/(25\pi^2)$ is a theorem, not a measurement.

%=====================================================================
\subsection{Implications}
%=====================================================================

\subsubsection{GUT symmetry breaking reinterpreted}

``Breaking'' is a misnomer. The chiral structure $(2,3)$ exists at all $N$. What changes with $N_\text{eff}$ is its spectral visibility. GUT breaking is inevitable (N$_\text{eff}$ always crosses $\sim$1500), mechanism-free, and irreversible.

\subsubsection{The $\det(G_h)$ contributions identified}

The $\sim$2.4\% $\det(G_h)$ geometric contributions throughout DRLT --- in the proton mass, water bond angle, baryon asymmetry, and coupling constants --- are now identified as contributions from $\lambda_6, \ldots, \lambda_{d_\text{indep}}$: spectrally indistinguishable from $\lambda_1, \ldots, \lambda_5$ but representation-theoretically trivial.

The trace conservation $\sum_i \Delta_i = 0$ follows from $\tr(G) = N$, which holds regardless of how many dimensions are trivial. The same $\agut$ contribution appears in:

\begin{center}
\begin{tabular}{lccl}
\hline
Quantity & Leading & $+\,\agut$ & Residual \\
\hline
$m_p$ & 924 MeV & 937.5 MeV & 0.08\% \\
$\theta_{\text{H}_2\text{O}}$ & $104.48^\circ$ & $104.52^\circ$ & $0.00^\circ$ \\
$\eta_B$ & $5.98 \times 10^{-10}$ & $6.13 \times 10^{-10}$ & 0.2\% \\
\hline
\end{tabular}
\end{center}

These corrections are the signature of a unified theory: they vanish when gravity (the $\det$ contribution) is included alongside quantum mechanics.

\subsubsection{No anthropic principle}

The constants \emph{cannot} be different. $\text{SU}(3)\times\text{SU}(2)\times\text{U}(1)$ is the only non-trivial chiral structure. Its coupling constant $6/(25\pi^2)$ is determined by $\CC$. ``Why is there something rather than nothing?'' reduces to: ``Do complex numbers exist?''
